\import{relation.tex}
\import{relation/properties.tex}

\section{Quasiorders}

Quasiorders are also called preorders.

\begin{abbreviation}\label{quasiorder}
    $R$ is a quasiorder iff
    $R$ is a relation and $R$ is quasireflexive and transitive.
\end{abbreviation}

\begin{classification}\label{quasiorder_is_relation}
    Let $R$ be a quasiorder.
    Then $R$ is a relation.
\end{classification}
\begin{proof}
\end{proof}

\begin{abbreviation}\label{quasiorder_on}
    $R$ is a quasiorder on $A$ iff
        $R$ is a binary relation on $A$ and
        $R$ is reflexive on $A$ and transitive.
\end{abbreviation}

\begin{proposition}\label{quasiorder_on_implies_quasiorder}
    Suppose $R$ is a quasiorder on $A$.
    Then $R$ is a quasiorder.
\end{proposition}
\begin{proof}
    $R$ is a binary relation on $A$ by definition.
    $R$ is a relation by
        \auto{relation,subseteq,times_elem_is_tuple}.
    Follows by
        \auto{quasireflexive,reflexive_on,subseteq,times_tuple_elim}.
\end{proof}

\begin{struct}\label{quasiordered_set}
    A quasiordered set $X$ is a onesorted structure equipped with
    \begin{enumerate}
        \item\label{quasiorder_relation}
            a nonstrict order relation
            $\leq[X]\subseteq\carrier[X]\times\carrier[X]$;
    \end{enumerate}
    such that
    \begin{enumerate}
        \item\label{quasiorder_refl}
            $\leq[X]$ is reflexive on $\carrier[X]$.
        \item\label{quasiorder_tran}
            $\leq[X]$ is transitive.
    \end{enumerate}
\end{struct}

\begin{classification}\label{quasiorder_relation_transitive}
    Let $X$ be a quasiordered set.
    Then $\leq[X]$ is transitive.
\end{classification}
\begin{proof}
    Follows by \obv{quasiorder_tran}.
\end{proof}

\begin{proposition}\label{quasiorder_type}
    Let $X$ be a quasiordered set.
    Then $\leq[X]$ is a binary relation on $\carrier[X]$.
\end{proposition}
\begin{proof}
    Follows by \obv{quasiorder_relation}.
\end{proof}

\begin{lemma}\label{quasiorder_transitive_double}
    Let $X$ be a quasiordered set.
    Let $a,b,c,d\in\carrier[X]$.
    Suppose $a\mathrel{\leq[X]}b\mathrel{\leq[X]}c
        \mathrel{\leq[X]}d$.
    Then $a\mathrel{\leq[X]} d$.
\end{lemma}
\begin{proof}
    $\leq[X]$ is transitive by \obv{quasiorder_tran}.
    Thus $a\mathrel{\leq[X]} c\mathrel{\leq[X]} d$ by \obv{transitive}.
    Hence $a\mathrel{\leq[X]} d$ by \obv{transitive}.
\end{proof}

\begin{proposition}\label{subseteqrel_is_quasiorder}
    $\subseteqrel{A}$ is a quasiorder on $A$.
\end{proposition}
\begin{proof}
    Show for all $w\in\subseteqrel{A}$ we have $w\in A\times A$.
    \begin{subproof}
        Fix $w\in\subseteqrel{A}$.
        Take $a,b$ such that $w=(a,b)$ and $a,b\in A$
            by \auto{subseteqrel,pair_eq_iff}.
        Follows by \auto{times_tuple_intro}.
    \end{subproof}
    Thus $\subseteqrel{A}$ is a binary relation on $A$
        by \obv{subseteq}.
    $\subseteqrel{A}$ is reflexive on $A$
        by \obv{reflexive_on,subseteqrel,subseteq}.
    Show for all $a,b,c$ such that
        $a\mathrel{\subseteqrel{A}}b$ and
        $b\mathrel{\subseteqrel{A}}c$
        we have $a\mathrel{\subseteqrel{A}}c$.
    \begin{subproof}
        Fix $a,b,c$.
        Assume $a\mathrel{\subseteqrel{A}}b$ and
            $b\mathrel{\subseteqrel{A}}c$.
        Take $x,y$ such that $(a,b)=(x,y)$ and
            $x,y\in A$ and $x\subseteq y$ by \auto{subseteqrel}.
        Then $a=x$ and $b=y$ by \obv{pair_eq_iff}.
        Thus $a\subseteq b$ and $a\in A$.
        Take $u,v$ such that $(b,c)=(u,v)$ and
            $u,v\in A$ and $u\subseteq v$ by \auto{subseteqrel}.
        Then $b=u$ and $c=v$ by \obv{pair_eq_iff}.
        Thus $b\subseteq c$ and $c\in A$.
        Thus $a\subseteq c$ by \obv{subseteq_transitive}.
        Follows by \obv{subseteqrel}.
    \end{subproof}
    Thus $\subseteqrel{A}$ is transitive
        by \obv{transitive}.
\end{proof}

\begin{proposition}\label{converse_quasiorder_on}
    Suppose $R$ is a quasiorder on $A$.
    Then $\converse{R}$ is a quasiorder on $A$.
\end{proposition}
\begin{proof}
    $R\subseteq A\times A$ and $R$ is reflexive on $A$ and
        $R$ is transitive by definition.
    $\converse{R}\subseteq A\times A$ by \obv{converse_type}.
    Hence $\converse{R}$ is a binary relation on $A$.
    $\converse{R}$ is reflexive on $A$ by
        \auto{reflexive_on,converse_iff}.
    Show that for all $x,y,z$ such that
        $x\mathrel{\converse{R}}y$ and
        $y\mathrel{\converse{R}}z$ we have
        $x\mathrel{\converse{R}}z$.
    \begin{subproof}
        Fix $x,y,z$.
        Assume $x\mathrel{\converse{R}}y$ and
            $y\mathrel{\converse{R}}z$.
        Hence $y\mathrel{R}x$ by \obv{converse_iff}.
        Hence $z\mathrel{R}y$ by \obv{converse_iff}.
        Hence $z\mathrel{R}x$ by \obv{transitive}.
        Follows by \obv{converse_iff}.
    \end{subproof}
    Hence $\converse{R}$ is transitive by \obv{transitive}.
\end{proof}
